\documentclass[12pt]{article}
\usepackage{jmlr2e}
\let\circledR\relax\let\circledS\relax\usepackage{newcomputermodern}
\usepackage{amsmath,amssymb,enumerate,booktabs,lastpage}

\newtheorem{hypothesis}{Hypothesis}

\jmlrheading{27}{2026}{1-\pageref{LastPage}}{1/26; Revised 5/26}{9/26}{26-0000}{Masum Billal}

\ShortHeadings{Global Entity Resolution}{Masum Billal}
\firstpageno{1}

\begin{document}
	
	\title{Global Coreference Resolution via Local Resolution and Quotient-Graph Inference}
	
	\author{\name Masum Billal \email billalmasum93@gmail.com \addr}
	
	\editor{TBD}
	
	\maketitle
	
	\begin{abstract}
		Modern coreference systems typically assume that document-level entity resolution requires global quadratic attention over the entire text. This paper challenges such assumptions by showing that most coreference decisions can be resolved locally and that long-range entity structure can be recovered through structured inference over compressed cluster representations. We introduce a two-stage inference framework that constructs a quotient space over local equivalence classes of mentions, followed by structured cross-window lifting of equivalence relations via cluster-level matching in a compressed graph representation. The model employs frozen semantic and contextual encoders throughout, training only small task-specific components. Evaluated on the CoNLL-2012 benchmark with gold mentions, the proposed approach achieves a CoNLL F$_1$ score of 0.8556 while substantially reducing memory usage and computational requirements relative to full-document architectures. These results suggest that entities can be resolved globally from localized reasoning and structured aggregation, without requiring global token-level interaction modeling over the full document, instead operating on a hierarchically constructed quotient graph of mention equivalence classes.
	\end{abstract}
	
	
	\begin{keywords}
		coreference resolution, representation factorization, quotient graph, frozen language models, window-based inference, transitive closure, efficient inference
	\end{keywords}
	
	\section{Introduction}\label{sec:intro}
	Coreference resolution is commonly defined as the task of clustering a set of entity mentions $\mathcal{M} = \{m_1, \ldots, m_n\}$ into groups that correspond to underlying real-world entities. Formally, this induces an unknown equivalence relation $\sim$ over mentions, where $m_i \sim m_j$ if and only if they refer to the same entity, and the output of the system is a partition $\mathcal{E} = \mathcal{M} / \sim$. From this perspective, coreference can be viewed as a structured prediction problem whose objective is to recover a latent quotient structure rather than to independently classify mention pairs or model token-level interactions.
	
	In this work, we explicitly adopt this equivalence-relation formulation as the primary object of inference and study how it can be recovered from localized computations over restricted subsets of $\mathcal{M}$. We propose viewing the problem as a two-level decomposition: (i) construction of partial equivalence relations under localized observation constraints, and (ii) recovery of the global partition via composition and closure over these local relations in a compressed representation space. This formulation separates local relational induction from global consistency enforcement, allowing long-range dependencies to be recovered through structured aggregation over induced equivalence classes rather than direct global interaction over the full mention set.
	
	The remainder of this paper is organized as follows. Section~\ref{sec:overview} outlines the overall workflow of the two-stage pipeline. Section~\ref{sec:hypothesis} develops the mathematical foundations of representation factorization and locality-transitivity equivalence. Sections~\ref{sec:stagea} and~\ref{sec:stageb} detail the technical implementation of the Stage A local inference procedure and Stage B quotient-space matching modules respectively. Section~\ref{sec:setup} describes experimental setup, and Section~\ref{sec:results} presents empirical results and analysis.
	
	To operationalize this formulation, we construct a structured inference procedure that separates local relational induction from global consistency resolution. The system does not rely on a single unified contextual model over the full document; instead, it defines a factorized representation space over which local equivalence relations are computed. Each mention $m_i \in \mathcal{M}$ is encoded using a fixed set of pretrained representation functions that provide complementary views of the same input (contextual and lexical), which are kept separate to preserve distinct informational signals. These representations are not updated during training and are treated as fixed feature mappings into a shared inference space.
	
	Importantly, we assume no access to global interaction over the full mention set. Instead, computation is restricted to bounded observation domains over which partial evidence of $\sim$ is induced. We introduce a set of bounded observation domains $\mathcal{W} = \{W_1, \ldots, W_K\}$ over $\mathcal{M}$. These domains do not define the equivalence structure itself; rather, they serve as local contexts in which restricted instantiations of the global relation are computed. Within each subset $W_k$, we estimate a local relation $\sim_k$ that approximates the restriction of $\sim$ to that region.
	
	Only after defining these local relations do we obtain structured intermediate objects that support global inference. Specifically, Stage A constructs local relational structure within each $W_k$, producing locally consistent equivalence classes that form a compressed intermediate representation of $\mathcal{M}$.
	
	After constructing local equivalence relations within each partition $W_k$, we obtain induced local classes that serve as the basis for global reasoning. These induced structures are not yet globally consistent, but they define a reduced representation space over which cross-region inference is performed.
	
	Stage B is therefore formulated as a quotient-space completion problem, where the goal is to recover global consistency of the induced equivalence relation by operating over the set of local equivalence classes rather than individual mentions. Instead of defining interactions over $\mathcal{M}$, we define a lifted interaction space over $\mathcal{C}$, in which each element corresponds to an entire local equivalence class. Relations between these elements are computed using aggregate compatibility over their constituent members, but are interpreted as relations between equivalence classes in the quotient space rather than as pairwise mention alignments.
	
	The resulting structure defines a higher-level relation $\approx$ over $\mathcal{C}$, which approximates the quotient of the unknown global equivalence relation induced on $\mathcal{M}$. Global coreference resolution is obtained by computing the transitive closure of $\approx$, yielding a consistent refinement of the initial local partitions. In this view, Stage B is not a post-hoc clustering step over mention groups, but a lifting operator that completes a partially specified equivalence relation in the quotient space induced by Stage A.
	
	\section{Solution Overview}\label{sec:overview}
	
	To bypass global quadratic text attention, the framework processes an incoming document through a two-stage hierarchical inference pipeline that operates first on bounded local context and then on a compressed quotient representation of induced entities.
	
	\begin{enumerate}
		\item \textbf{Decoupled Encoding Stream:} 
		Each mention is mapped into two fixed, pretrained representation spaces capturing complementary semantic and contextual information. These encoders are frozen throughout training, and their outputs are never fused into a single joint representation at the token level. All subsequent reasoning is performed over these fixed feature views.
		
		\item \textbf{Stage A (Intra-Window Relational Construction):} 
		The document is partitioned into a set of non-overlapping windows of fixed size. Within each window, all pairwise mention interactions are computed using a local scoring function over the frozen representations.
		
		A transitive inference procedure (union-find) is then applied strictly within each window, producing a set of locally consistent equivalence structures. These structures define the first level of inferred entity groupings and constitute the initial compressed representation of the document.
		
		\item \textbf{Quotient Representation Induction:} 
		The output of Stage A is a collection of locally inferred equivalence groups, each representing a collapsed subset of mentions. These groups are treated as atomic units for all subsequent computation, inducing a reduced representation space in which reasoning operates over clusters rather than raw mentions.
		
		\item \textbf{Stage B (Cross-Window Quotient Matching):} 
		Global entity consistency is recovered by evaluating interactions between the induced clusters rather than individual mentions. A size-invariant compatibility function computes relational evidence between clusters using aggregated mention-level interactions combined with auxiliary lexical signals.
		
		Edges exceeding a learned threshold define cross-window equivalences between clusters. Final document-level entities are obtained by computing the transitive closure over this cluster-level graph, yielding a globally consistent partition of mentions.
	\end{enumerate}
	\section{Theoretical Foundation and Design Principles}\label{sec:hypothesis}
	
	\subsection{The Representation Factorization Hypothesis}
	
	Monolithic language models rely on end-to-end contextualization over full-document sequences to implicitly disentangle semantic, discourse, and referential dependencies. This couples representation learning with relational inference and induces unnecessary computational cost. We instead hypothesize that coreference can be reformulated as a structured inference problem over localized observations combined through a compressed global reasoning space. Under this view, global entity structure does not require direct token-level interaction, but can be recovered from compositions of locally induced relational evidence. This motivates a factorization of inference into independent information channels that preserve complementary signals for downstream relational reconstruction.
	
	\begin{hypothesis}[Low-Redundancy Information Basis Hypothesis]
		Let $\mathcal{V}$ denote the decision space required to recover global entity structure. We assume it can be decomposed into $d$ partially independent feature channels:
		\[
		\mathcal{V} = \mathcal{V}_1 \times \mathcal{V}_2 \times \dots \times \mathcal{V}_d.
		\]
		Each channel is assumed to contain non-redundant information useful for relational reconstruction.
	\end{hypothesis}
	
	Independence here refers to functional non-redundancy rather than orthogonality: features are valuable insofar as they cannot be reconstructed from other channels without loss of task-relevant information.
	
	Empirically, this is supported by a diagnostic comparison in which scalar similarity over embedding axes yields weak separability (AUC ≈ 0.57), while a learned linear probe over separated representation channels achieves substantially higher performance (AUC = 0.91).
	
	\subsection{Locality-Transitivity Equivalence}
	
	We further hypothesize that long-range coreference does not require direct interaction modeling, but can be reconstructed via transitive propagation over locally observable relational chains.
	
	\begin{hypothesis}[Transitive Link Reconstruction]
		Global entity identity can be decomposed into chains of locally observable mention-level relations such that any two coreferent mentions are connected via a finite sequence of intermediate mentions, each co-located within some shared observation context.
	\end{hypothesis}
	
	\begin{theorem}[Quotient Reconstruction Principle]
		If local relational structure is correctly identified within bounded contexts and cross-context correspondences are correctly recovered, then global entity structure can be reconstructed exactly via transitive closure over the induced relational graph.
	\end{theorem}
	
	This implies that global quadratic attention is not strictly required for entity recovery, but can be replaced by structured inference over localized relational graphs.
	\section{Stage A: Local Intra-Window Resolution}
	\label{sec:stagea}
	Let a document be represented as a set of mentions
	\[
	\mathcal{M} = \{m_1, \dots, m_N\}.
	\]
	
	We define a deterministic partition into disjoint observation windows:
	\[
	\mathcal{W} = \{W_1, \dots, W_K\}, \quad
	W_i \cap W_j = \emptyset, \quad
	\bigcup_{k=1}^{K} W_k = \mathcal{M}.
	\]
	
	Each mention is assigned to exactly one window via:
	\[
	\pi: \mathcal{M} \rightarrow \mathcal{W}.
	\]
	
	Each mention is represented using a fixed factorized encoder:
	\[
	\mathbf{h}_i = \left[ f_{\theta_1}(m_i) \parallel f_{\theta_2}(m_i) \right],
	\]
	where both encoders are frozen.
	
	Within each window $W_k$, we define a local relation:
	\[
	m_i \sim_k m_j \iff s_\theta(\mathbf{h}_i, \mathbf{h}_j) > \tau.
	\]
	
	Let $\sim_k^*$ denote its transitive closure, inducing local clusters:
	\[
	\mathcal{C}_k = W_k / \sim_k^*.
	\]
	
	We define the full set of local clusters:
	\[
	\mathcal{C} = \bigcup_{k=1}^{K} \mathcal{C}_k.
	\]
	
	Each cluster $C \in \mathcal{C}$ is treated as an atomic unit for subsequent inference. Stage A computes the local equivalence structure defined in Section~\ref{sec:model_class}. Each window $W_k \in \mathcal{W}$ is processed independently. For each pair of mentions $(m_i, m_j) \in W_k \times W_k$, a compatibility score is computed:
	\[
	s_\theta(\mathbf{h}_i, \mathbf{h}_j).
	\]
	
	A local relation is induced via thresholding:
	\[
	m_i \sim_k m_j \iff s_\theta(\mathbf{h}_i, \mathbf{h}_j) > \tau.
	\]
	
	Transitive closure over $\sim_k$ yields:
	\[
	\mathcal{C}_k = W_k / \sim_k^*.
	\]
	
	This induces the set of local clusters:
	\[
	\mathcal{C} = \bigcup_{k=1}^{K} \mathcal{C}_k,
	\quad \mathcal{Q} \equiv \mathcal{C}.
	\]
	
	These clusters form the atomic units used in Stage B.
	
	\section{Stage B: Cross-Window Quotient Matching}\label{sec:stageb}
	Stage B reconciles entity fragments across window pairs, operating entirely over the quotient graph nodes produced by Stage A. It scores cluster pairs directly over their full cross-cluster mention grid via an architecture optimized for fine-grained interaction (\texttt{MentionMatcher}).
	
	To enable consistent reasoning over the quotient space, we define a lifted compatibility operator over local equivalence classes that preserves the structure induced by mention-level relations. Let $A, B \in \mathcal{C}$ denote two equivalence classes produced by Stage A. Since each class aggregates multiple mentions, relations between $A$ and $B$ must be defined in a way that is invariant to internal class cardinality and robust to heterogeneous mention composition.
	
	We define a quotient-compatible interaction operator $\sigma : \mathcal{C} \times \mathcal{C} \rightarrow \mathbb{R}$ that aggregates fine-grained mention-level interactions into a single class-level relational quantity while respecting permutation invariance within each class. Rather than collapsing each class into a single vector representation, we define the interaction directly over the full bipartite structure induced by their constituent mentions:
		\begin{align*}
			\sigma(A,B)
				& = \log\sum_{m\in A}\sum_{n\in B}\exp\left(f(m,n)\right)-\log(|A|\cdot|B|)
		\end{align*}
	where $f(m,n)$ is a local compatibility functional defined over the fixed representations of mentions. This formulation can be interpreted as a log-partition lifting operator that maps pairwise mention interactions into a stable class-level relation in the quotient space, while explicitly normalizing for class size to avoid structural bias toward larger equivalence classes.
	
	In this view, $\sigma(A,B)$ is not a similarity score between embeddings, but a measure of relational support between equivalence classes induced by aggregated evidence over their member mentions. The use of a log-sum-exp aggregation ensures that strong local alignments between any pair of mentions can dominate the class-level relation when supported consistently, while still allowing distributed weaker evidence to contribute additively. This defines a stable lifting of mention-level relations into the quotient space $\mathcal{C}$, enabling global equivalence completion via transitive closure over class-level interactions.
	\subsection{Lexical Features}
	To assist the frozen contextual spaces across deep token spans, the interaction head is augmented with an additive combination of localized, discrete string-identity features:
	\begin{align*}
		\sigma(A, B) = \text{score}_{\text{neural}}(A, B) + \mathbf{w}^\top \mathbf{x}_{\text{lex}}
	\end{align*}
	where $\mathbf{w}$ is a learned weight vector, and $\mathbf{x}_{\text{lex}}$ contains:
	\begin{enumerate}
		\item \textbf{IDF-weighted token Jaccard similarity:} $\sum_{t \in A \cap B} \text{idf}(t) / \sum_{t \in A \cup B} \text{idf}(t)$, where document-local $\text{idf}(t) = \log(M / (1 + \text{df}(t))) + 1$.
		\item \textbf{Substring containment:} A binary indicator marking if any content mention (excluding bare pronouns) in $A$ is a substring of any content mention in $B$, or vice versa.
		\item \textbf{Exact mention match:} A binary flag tracking if the two clusters share an identical content mention surface text string.
	\end{enumerate}
	Final document clusters are constructed via global union-find over all cluster pairs whose final score $\sigma(A, B)$ exceeds a learned null bias parameter.
	
	\section{Experimental Setup}\label{sec:setup}
	All evaluation experiments are conducted on the test split of the CoNLL-2012 English coreference benchmark \citep{pradhan2012conll} utilizing gold mention boundaries. Following empirical data mixture optimizations, our lightweight scoring modules are trained on a uniform corpus mixture consisting of CoNLL-2012, LitBank, CorefUD, and a subsampled partition of 8,000 documents from PreCo. Training enforces a uniform loss across all datasets without negative window-pair subsampling to preserve distributional balance. All baseline results cited from prior literature correspond strictly to their respective gold-mention evaluation configurations.
	
	\section{Results}\label{sec:results}
	
	\subsection{Main Coreference Benchmarks}
	Table \ref{tab:stageb} establishes the primary system performance evaluated against state-of-the-art architectures using identical gold mention contexts.
	
	\begin{table}[ht]
		\centering
		\small
		\begin{tabular}{lcccc}
			\toprule
			\textbf{Approach} & \textbf{CoNLL} & \textbf{MUC} & \textbf{B$^3$} & \textbf{CEAFe} \\
			\midrule
			Stage A Only (No Merge) & 0.7439 & 0.8460 & 0.6910 & 0.6950 \\
			\textbf{Stage A + B (Mention-Level)} & \textbf{0.8556} & \textbf{0.9230} & \textbf{0.8310} & \textbf{0.8140} \\
			\midrule
			\citet{dobrovolskii2021word}        & 0.8180 & 0.8570 & 0.8000 & 0.7970 \\
			\citet{caciularu2023peek}           & 0.8360 & 0.8740 & 0.8200 & 0.8130 \\
			\citet{bohnet2023coreference}       & 0.8380 & 0.8760 & 0.8220 & 0.8150 \\
			\citet{luo2025imcoref}              & 0.8510 & 0.8890 & \textbf{0.8350} & \textbf{0.8290} \\
			\bottomrule
		\end{tabular}
		\caption{Full-document CoNLL-2012 test results on \textbf{gold mention boundaries}. Baseline systems are mapped strictly to their gold-mention settings to isolate linking capability.}
		\label{tab:stageb}
	\end{table}
	
	Table \ref{tab:stageb_type} itemizes performance across granular cluster configurations, tracking the architectural stability of the model over varied mention types.
	
	\begin{table}[ht]
		\centering
		\small
		\begin{tabular}{lcccc}
			\toprule
			\textbf{Cluster Category Type} & \textbf{CoNLL} & \textbf{MUC} & \textbf{B$^3$} & \textbf{CEAFe} \\
			\midrule
			PROPN-only & 0.8390 & 0.8400 & 0.8380 & 0.8390 \\
			PRON+NOUN  & 0.7700 & 0.7630 & 0.7490 & 0.7970 \\
			NOUN-only  & 0.7500 & 0.7530 & 0.7500 & 0.7480 \\
			PRON+PROPN & 0.7280 & 0.7450 & 0.6940 & 0.7430 \\
			all-mixed  & 0.7160 & 0.7730 & 0.6730 & 0.7030 \\
			NOUN+PROPN & 0.7130 & 0.7190 & 0.7040 & 0.7150 \\
			PRON-only  & 0.6640 & 0.7050 & 0.6490 & 0.6370 \\
			\bottomrule
		\end{tabular}
		\caption{Definitive per-cluster category breakdown for the Stage A + B pipeline, CoNLL-2012 test.}
		\label{tab:stageb_type}
	\end{table}
	\subsection{Complexity Analysis}
	Since window size is fixed, $W=N/K$. Assuming mentions are approximately uniformly distributed, $m_{w}\approx M/W=MK/N$. Then total encoding complexity is
		\begin{align*}
			O\left(W*K^{2}\right)
				& = O(NK)
		\end{align*}
	instead of $O(N^{2})$. in stage A, computation cost is
		\begin{align*}
			O\left(w\left(\dfrac{M}{w}\right)^{2}\right)
				& = O\left(\dfrac{M^{2}K}{N}\right)
		\end{align*}
	This is essentially linear until stage A. For stage B, complexity is
		\begin{align*}
			O(C^{2})
		\end{align*}
	Total complexity becomes
		\begin{align*}
			O\left(NK+\dfrac{M^{2}K}{N}+C^{2}\right)
		\end{align*}
	For a typical document, we realistically have $C\ll M\ll N$. Therefore, even though quadratic computation is not fully eliminated, we only do it in a much more reduced space. Mean document statistics on CoNLL-2012: $N = 617.4$ tokens, $M = 57.1$ mentions $C = 20.7$ quotient nodes. The quotient graph contains only $20.7 / 57.1 = 36.3\%$ of the original mention nodes. Relative to token-level reasoning, the reduction is $617.4 / 20.7 = 29.8$ times per dimension, corresponding to roughly $(617.4 / 20.7)^2 ≈ 890$ times fewer pairwise comparisons.
	\begin{table}[ht]
		\centering
		\small
		\begin{tabular}{lcc}
			\toprule
			Space & Nodes & Pairwise Comparisons \\
			\midrule
			Token Space & 617.4 & 381,183 \\
			Mention Space & 57.1 & 3,260 \\
			Quotient Space & 20.7 & 428 \\
			\bottomrule
		\end{tabular}
		\caption{Average document graph sizes on CoNLL-2012.}
	\end{table}
	The empirical results validate this generalized theoretical framework and information routing properties:
	\begin{enumerate}
		\item \textbf{Validation of Independent Basis Projections (Hypothesis 1):} Our core framework reaching a definitive $0.8556$ $F_1$ score demonstrates that deep contextual models do not need to fine-tune internal weights to parse cross-cutting signals. By treating specialized universal embeddings as structural axes and preserving separate vector paths, the linear combination layer tracks identity without collapsing coordinates. This is strongly substantiated by our linear probe diagnostic, where keeping distinct vector coordinates yields an AUC of $0.91$, compared to the near-chance level of $\approx 0.57$ encountered when collapsing dimensions to scalar cosine steps.
		\item \textbf{Validation of Locality (Hypothesis 2 \& Proposition 1):} Stage A alone (operating strictly inside isolated windows without a global merging layer) yields a baseline score of $0.7439$ F$_1$. This means that local-only tracking recovers approximately 87\% ($0.7439 / 0.8556$) of the complete global framework performance. This explicitly confirms that entity dependencies are highly concentrated within bounded token ranges, proving that local quadratic attention is sufficient for core resolution.
		\item \textbf{Validation of Quotient Graph Efficiency:} Evaluating the structure of the document space shows that Stage A collapses an average of $57.1$ mentions per document into just $20.7$ cluster nodes. This produces a structural compression profile of \textbf{2.76:1} before cross-window matching starts.
		\item \textbf{Validation of Transitive Linking:} Stage B's matching head resolves cross-window splits with an Edge Precision of 79.79\% ($2567 / 3217$) and an Edge Recall of 70.96\% ($2883 / 4063$). Each true entity fragments into only $1.45$ components on average, demonstrating that transitive matching over compressed quotient nodes can successfully reconstruct long-range entities.
	\end{enumerate}
	
	\subsection{Resource Footprint and Efficiency Analysis}
	By locking the primary encoder backbones and routing cross-window connections through the quotient graph, the model significantly reduces computational overhead. 
	
	\begin{table}[ht]
		\centering
		\small
		\begin{tabular}{lrc}
			\toprule
			\textbf{Component} & \textbf{Parameter Count} & \textbf{Status} \\
			\midrule
			Stage A (AntecedentScorer) & 8.43M & Trained \\
			Stage B (MentionMatcher)    & 8.40M & Trained \\
			RoBERTa-large               & 355.00M & Frozen \\
			BGE-large                   & 335.00M & Frozen \\
			\midrule
			\textbf{Total Parameters}   & \textbf{706.83M} & (2.38\% Trained) \\
			\bottomrule
		\end{tabular}
		\caption{System parameter budget distribution.}
		\label{tab:parameters}
	\end{table}
	
	As itemized in Table \textbf{\ref{tab:parameters}}, trained parameters account for only 2.38\% (16.8M) of the total parameter budget. In inference, this architectural factorization allows the model to achieve a wall-clock throughput of 19.4 ms / 1k tokens on consumer laptop hardware (a single RTX 4060 Laptop GPU). Crucially, the windowed compression layers prevent the system from materializing the full mention-pair score matrix, lowering peak inference VRAM requirements by \textbf{3.1$\times$} compared to full-document all-pairs execution (159 MB vs 488 MB).
	
	\section{Limitations and Future Work}\label{sec:future}
	Every result in this paper leaves both underlying representation encoders frozen. Consequently, the contextual channel is entirely dependent on general-purpose RoBERTa-large activations with zero coreference-specific fine-tuning. The PRON-only floor ($0.6640$) highlights this boundary: pronouns carry low semantic signal, forcing the matching head to rely on contextual binding spaces that naturally decay across extreme window gaps. Furthermore, this work assumes gold mention boundaries to isolate coreference resolution logic from mention detection errors. Scaling this quotient-graph architecture to consume predicted mentions from an open span detector remains an immediate avenue for future research.
	
	\section{Conclusion}\label{sec:conclusion}
	This work demonstrates that global entity resolution can be recovered from local inference plus quotient-graph reconstruction. More importantly, the results provide evidence for a broader computational hypothesis: many language tasks may be decomposable into low-redundancy information channels whose interactions can be modeled independently and recombined through lightweight composition layers. Under such a decomposition, expensive global attention need not operate over the original token space, but only over a compressed latent state space induced by the task structure.
	
	\bibliography{references}
	
\end{document}